\section*{Abstract}
Do the signal-engineering representations (Doc.~288) help to better capture quarks and neutrinos? The
honest answer: \emph{as language and diagnosis yes, as derivation no} -- and very differently for the
two sectors. Each fermion triplet gets a magnitude coordinate (Koide $Q$, i.e.\ $r=\sqrt{6Q-2}$) and a
phase coordinate. Computed: only the charged leptons sit at the special point $r=\sqrt2$ ($Q=2/3$); the
quarks lie elsewhere ($r\approx1{.}54$ resp.\ $1{.}76$, and scale-dependent), and for neutrinos $Q$
cannot even be formed. The gain is a \emph{unifying map}: leptons and quarks are magnitude-dominated
(large hierarchy, small mixing), neutrinos phase-dominated (flat/degenerate magnitude, large mixing --
PMNS/CKM weight $0{.}55$ vs $0{.}035$, factor~16). FFGFT's own neutrino hypothesis (equal masses +
oscillation from a geometric phase) is then literally the \emph{pure all-pass sector}: magnitude flat,
all physics in the phase -- and it becomes falsifiable, because the geometric phase must reproduce the
measured $\Delta m^2$ phases.

\section{The question}
Doc.~288 showed that the Z$_3$-circulant mass operator gives concrete computational levers via the DFT.
The natural question: do the same representations reach the hard sectors -- quarks and neutrinos? An
honest stock-taking is needed first, because the corpus does not treat these sectors like the clean
charged leptons.

\section{What the corpus already does (per sector)}
\begin{itemize}
  \item \textbf{Charged leptons} (Doc.~006/046/258): the clean case -- Foot-Koide form, $Q=2/3$,
        $r=\sqrt2$, Z$_3$-circulant. Everything from Doc.~288 applies directly.
  \item \textbf{Quarks} (Doc.~006/005): the same $r_i/p_i$ geometry in principle, but the direct
        geometric method reaches only $\sim1{.}2\%$; for quarks it needs the \emph{extended fractal
        method} with QCD dynamics and \emph{ML calibration on lattice data} (FLAG~2024). Hence not
        purely geometric -- and quark masses \emph{run} (scale/scheme dependent).
  \item \textbf{Neutrinos} (Doc.~007/047): a \emph{different} mechanism -- photon analogy,
        $m_\nu=\tfrac{\xi_0^2}{2}m_e\approx4{.}54$~meV, all three flavours (nearly) equal in mass,
        oscillation from \emph{geometric phases} rather than mass differences. Flagged throughout the
        corpus as \textbf{speculative, without empirical basis}.
\end{itemize}

\section{The (Q, r) coordinate: where the sectors sit}
Via $Q=\tfrac13+\tfrac{r^2}{6}$ (so $r=\sqrt{6Q-2}$) every triplet gets a magnitude coordinate.
Computed (\texttt{koide\_q\_sektoren.py}, PDG-near values):

\begin{table}[htbp]
\centering
\caption{Magnitude coordinate of the three formable triplets. Quark values are scale-dependent (MS-bar).}
\label{tab:qr}
\begin{tabular}{p{5.2cm}p{2.2cm}p{2.4cm}p{2.6cm}}
\toprule
\textbf{Sector} & \textbf{Koide $Q$} & \textbf{$r=\sqrt{6Q-2}$} & \textbf{Position} \\
\midrule
charged leptons (e,\,$\mu$,\,$\tau$) & $0{.}6667$ & $\sqrt2=1{.}4142$ & special \\
up quarks (u,\,c,\,t)                & $0{.}849$  & $1{.}759$         & not special \\
down quarks (d,\,s,\,b)              & $0{.}730$  & $1{.}543$         & not special \\
neutrinos                            & ---        & ---               & no triplet \\
\bottomrule
\end{tabular}
\end{table}

This is already the first answer: \emph{only the charged leptons sit at the special point}
$r=\sqrt2$. The quark $r$ are not recognisably special (and run with the scale); for neutrinos the
absolute scale and ordering are missing, so $Q$ cannot be formed -- only $\Delta m^2_{21}\approx
7{.}5\cdot10^{-5}$ and $\Delta m^2_{31}\approx2{.}5\cdot10^{-3}~\mathrm{eV}^2$ are measured.

\section{The magnitude/phase map}
The signal-engineering decomposition (magnitude $=$ mass spectrum, phase $=$ mixing/oscillation) places
the four sectors on \emph{one} map:
\begin{itemize}
  \item \textbf{Quarks}: strong mass hierarchy (large magnitude), small mixing (CKM near identity)
        -- \emph{magnitude-dominated}.
  \item \textbf{Charged leptons}: magnitude at a special value ($r=\sqrt2$), small phase
        ($\theta=2/9$) -- \emph{magnitude-dominated}.
  \item \textbf{Neutrinos}: flat magnitude (degenerate), large mixing (PMNS far from identity)
        -- \emph{phase-dominated}.
\end{itemize}
The CKM/PMNS contrast can be quantified (\texttt{ckm\_pmns\_kontrast.py}, off-diagonal weight of the
$|U|^2$ matrix): CKM $0{.}035$ vs PMNS $0{.}554$ -- a factor $\sim16$. The \enquote{small vs large
mixing} is thus exactly a \emph{magnitude vs phase} contrast.

\section{Neutrinos as the all-pass sector}
Here is the strongest point. FFGFT's own hypothesis -- equal masses, oscillation from a geometric phase
-- is, in magnitude/phase language, literally an \emph{all-pass}: $|H|=1$ (flat magnitude spectrum),
nontrivial phase. A flat magnitude spectrum cannot encode any oscillation; all physics sits in the
phase (\texttt{neutrino\_allpass.py}).

This makes the test \emph{concrete and falsifiable}. The standard oscillation phase is
$\Delta\phi_{ij}=1{.}267\,\Delta m^2_{ij}\,L/E$; at T2K-like $L/E$ ($L=295$~km, $E=0{.}6$~GeV) this is
$\Delta\phi_{31}\approx1{.}557$~rad and $\Delta\phi_{21}\approx0{.}047$~rad. FFGFT's geometric phase
(from $\tilde T\cdot m=1$, quantum numbers $n,\ell,j$) \emph{must} reproduce exactly these
$L/E$-dependent phases. If it does not, the degeneracy hypothesis is refuted.

\section{What it achieves -- and what it does not}
\begin{important}
\textbf{Achievement (diagnostic).} The magnitude/phase decomposition explains the \emph{division of
labour} between the sectors: why the leptons are clean (magnitude at a special value, small phase), the
quarks are not (magnitude not special, extreme hierarchy), the neutrinos quite different (magnitude
flat, all phase). This is genuine organising knowledge -- and it gives the speculative neutrino
hypothesis a sharp, falsifiable form.

\textbf{Boundary (P35).} This is better language, a unifying map and a sharper test -- \emph{not} a new
derivation of quark or neutrino masses. The quark $r$ are not special and scale-dependent, and the
quark sector in the corpus relies on QCD plus lattice calibration, not pure geometry. The neutrino
sector remains speculative (own flag, Doc.~007) and in tension with the very well-tested $\Delta m^2$
interpretation of oscillation. The clean all-pass map is a compelling reformulation, not new evidence.
\end{important}

\begin{conclusionBox}[Balance: one map, no new sector]
The signal-engineering representations capture quarks and neutrinos \emph{better} -- but in a precise
sense: they give a common magnitude/phase coordinate in which only the charged leptons sit at the
special point $r=\sqrt2$, the quarks magnitude-dominated beside them, and the neutrinos as a pure
phase / all-pass sector. The CKM/PMNS contrast becomes magnitude vs phase (factor~16). And FFGFT's own
neutrino hypothesis acquires a falsifiable form: the geometric phase must hit the measured $\Delta m^2$
oscillation phases. This is diagnosis and test, not a mass derivation -- honestly bounded.
\end{conclusionBox}
